% theme: meaning_and_notation_of_vectors
\MajorQuestionLesson{ベクトルの意味と表し方}
\SetRowSpace{4mm}

\TypeQuestionSingle{次の有向線分について、指定されたものを答えなさい。}{directed_segment}
\QQRowThree{$\overrightarrow{\mathrm{AB}}$ の始点}{}{$\overrightarrow{\mathrm{PQ}}$ の終点}{}{$\overrightarrow{\mathrm{XY}}$ の始点}{}
\QQRowThree{$\overrightarrow{\mathrm{CD}}$ の終点}{}{$\overrightarrow{\mathrm{RS}}$ の始点}{}{$\overrightarrow{\mathrm{UV}}$ の終点}{}

\TypeQuestionSingle{次のベクトルの大きさを表す記号を書きなさい。}{magnitude_notation}
\QQRowThree{$\vec{a}$}{}{$\vec{p}$}{}{$\overrightarrow{\mathrm{AB}}$}{}
\QQRowThree{$\vec{x}$}{}{$\overrightarrow{\mathrm{PQ}}$}{}{$-\vec{b}$}{}

\TypeQuestionDouble{次の量が、向きと大きさをもつ量なら「ベクトル」、大きさだけをもつ量なら「スカラー」と答えなさい。}{vector_or_scalar}
\QQRowThree{変位}{}{気温}{}{速度}{}
\QQRowThree{質量}{}{力}{}{時間}{}
\QQRowThree{加速度}{}{面積}{}{位置の変化}{}

\TypeQuestionSingle{次の文を、ベクトルの記号を用いて表しなさい。}{write_vector}
\QQRow{点 $\mathrm{A}$ から点 $\mathrm{B}$ へのベクトル}{}{点 $\mathrm{Q}$ から点 $\mathrm{P}$ へのベクトル}{}
\QQRow{点 $\mathrm{X}$ から点 $\mathrm{Y}$ への変位}{}{点 $\mathrm{D}$ から点 $\mathrm{C}$ へのベクトルの大きさ}{}

\LessonPageBreak
\SetRowSpace{4mm}

\TypeQuestionDouble{一直線上に点 $\mathrm{A},\mathrm{B},\mathrm{C},\mathrm{D}$ がこの順に等間隔で並んでいる。次のベクトルを $\overrightarrow{\mathrm{AB}}$ を用いて表しなさい。}{equally_spaced_points}
\QQRowThree{$\overrightarrow{\mathrm{BC}}$}{}{$\overrightarrow{\mathrm{CB}}$}{}{$\overrightarrow{\mathrm{AC}}$}{}
\QQRowThree{$\overrightarrow{\mathrm{DA}}$}{}{$\overrightarrow{\mathrm{BD}}$}{}{$\overrightarrow{\mathrm{DC}}$}{}

\TypeQuestionDouble{四角形 $\mathrm{ABCD}$ は平行四辺形である。次のベクトルと等しいベクトルを、頂点を用いて1つ答えなさい。}{parallelogram_vectors}
\QQRowThree{$\overrightarrow{\mathrm{AB}}$}{}{$\overrightarrow{\mathrm{AD}}$}{}{$\overrightarrow{\mathrm{BA}}$}{}
\QQRowThree{$\overrightarrow{\mathrm{BC}}$}{}{$\overrightarrow{\mathrm{CD}}$}{}{$\overrightarrow{\mathrm{DA}}$}{}

\TypeQuestionSingle{次のベクトルを、$\vec{a}$ または $-\vec{a}$ を用いて表しなさい。}{opposite_vector}
\QQRowThree{$\vec{a}=\overrightarrow{\mathrm{AB}}$ のとき $\overrightarrow{\mathrm{BA}}$}{}{$\vec{a}=\overrightarrow{\mathrm{PQ}}$ のとき $\overrightarrow{\mathrm{QP}}$}{}{$-\overrightarrow{\mathrm{CD}}=\vec{a}$ のとき $\overrightarrow{\mathrm{CD}}$}{}
\QQRow{$\overrightarrow{\mathrm{XY}}=-\vec{a}$ のとき $\overrightarrow{\mathrm{YX}}$}{}{$\vec{a}=-\overrightarrow{\mathrm{RS}}$ のとき $\overrightarrow{\mathrm{SR}}$}{}

\TypeQuestionDouble{次の記述が正しければ○、正しくなければ×を答えなさい。}{meaning_judgment}
\QQRow{大きさが等しい2つのベクトルは、必ず等しい。}{}{向きが同じ2つのベクトルは、必ず等しい。}{}
\QQRow{等しいベクトルは、始点が異なっていてもよい。}{}{零ベクトルの大きさは $0$ である。}{}
\QQRow{ベクトルとその逆ベクトルの大きさは等しい。}{}{零ベクトルには特定の向きを定めない。}{}
